\begin{center}
{\zihao{2}\bfseries 考虑变物性与径向收缩的药材\\[5pt]热风干燥建模及数值分析}
\end{center}
\vspace{3pt}
\begin{center}\bfseries 摘\quad 要\end{center}

中药是我国的文化瑰宝，其干燥工艺直接关系到品质与药效。该工艺的参数选取极大地影响干燥效率、能耗与成品品质。传统试验优化模式存在成本高、周期长的局限。为此，本文基于烘房的温湿环境与药材的尺寸变化规律，建立了在预热阶段和恒温阶段中，药材在圆柱坐标下一维轴对称的传热–传质模型，为探究药材的干燥规律提供了理论依据。

\textbf{针对问题一，建立预热平衡阶段的传热与水分扩散模型。}首先，由附录2得知热物性取常数、扩散系数只依赖含水率，二者互不耦合，故可分别求解；然后，以附件1的烘房温湿数据作分段线性插值驱动边界，列出药材表面对流换热与浓度差传质两个公式；之后，采用以半离散方程组和表面二次加密网格为核心的有限体积法；最终以变阶后向差分公式做隐式自适应时间积分，得到 1800s 内距药材中心不同位置的温度和水分浓度，结果保存到 result1.xlsx 中。

\textbf{针对问题二，建立变物性的热质耦合模型。}首先，根据附件1，拟合出恒温干燥阶段烘房的温度为50$^\circ$C和含水量为0.05kg/kg，该数据后两问通用；然后，将附录3中经验密度、比热容、导热系数和扩散系数嵌入局部传递方程；最终，通过界面通量和稀疏雅可比矩阵实现温湿状态的隐式同步更新，得到在恒温干燥阶段，各个时间点距药材中心不同位置的温度和水分浓度，并将结果保存在文件result2.xlsx中。

\textbf{针对问题三，建立基于变物性热质耦合的“干燥终点”模型。}本质是上一问的模型在时间维度上的延长。首先，判断“干燥终点”位置为圆柱中心轴线，并予以证明；其次，引入时间维度上的环境边界闭合假设；然后，建立全域最大含水率的阈值事件数学判据；最后，采用连续事件定位与空间插值实现数值求解与结果提取，得到干燥停止的时间为57.4740 h。

\textbf{针对问题四，建立药材的移动边界模型。}首先，引入移动边界是径向均匀收缩假设；其次，通过坐标变换推导材料坐标系下的控制方程与表面缩放边界，将移动域转化为固定计算域；然后，建立基于固定域坐标的全域最大含水率阈值事件判据；最后，采用固定网格下的自适应BDF求解与空间逆映射，得到药材在尺寸变化条件下的干燥停止时间为51.0920 h。

\vspace{5pt}
\noindent\textbf{关键词：}药材干燥；有限体积法；热质耦合；移动边界；自适应BDF

\nextpage
\section{问题背景与重述}
\subsection{问题背景}
作为中药材加工中至关重要的一道工序，热风干燥是个难题，烘房温湿条件选取的多样性赋予了"热风干燥"极高的工艺难度和探索价值。干燥时，热风从药材表面流过，热量由表及里传入内部，水分再由内向外迁移并在表面蒸发。在此期间，需要保持药材内部水分迁移通路的畅通，以得到最佳的干燥效果。

工艺参数选取不当（如升温速率与排湿节奏不匹配）常导致“外干内湿”、“表面结壳硬化”、能耗高昂及批次品质不稳定等工程难题。鉴于实体试验优化周期长、能耗大且难以直接观测柱体内部物性演化，建立基于非稳态传热传质偏微分方程的数理仿真模型，系统探究多变烘房环境下药材内部温湿度分布规律、动态烘干周期以及尺寸收缩效应，是实现中药材干燥从“经验调控”迈向“精准数理调控”的必然途径。

\subsection{问题重述}
中药材形状大致呈圆柱形，长为25 cm，半径为2 cm。烘干开始时，药材的温度为28°C ，水分浓度（即干基含水率）为2.55 kg/kg。

\textbf{问题一：}采用题目附录2的常热物性及含水率相关扩散系数，预测前1800 s的温度、干基含水率分布；在指定七个时刻与五个半径位置报告结果，并提供每隔1 s、0.1 cm的完整输出。

\textbf{问题二：}统一采用题目附录3的经验关系，建立温度与含水率相互作用的耦合模型；报告前3 h中每隔0.5 h五个径向位置的温度与含水率结果，并形成每隔1 s、0.1 cm的径向分布记录。

\textbf{问题三：}沿用问题二模型，以材料各处含水率均低于0.15 kg/kg为目标，确定干燥所需时间；按6 h间隔列出指定位置含水率，并保存每60 s、每0.1 cm的明细。

\textbf{问题四：}结合附件2中的半径变化及题目附录4的经验公式，建立药材的尺寸变化模型，确定烘干时长，报告每隔6小时药材不同径向位置与当前表面的水分分布，并保存每60 s、每0.1 cm的明细。

\section{问题分析}
\subsection*{问题一}
问题一限定于药材入烘房后的前 1800\,s，物性按附录 2 取常数，几何尺寸不变，要求给出该时段内温度与水分浓度沿半径方向的分布及其时间演化，并按规定的时刻与径向位置整理成结果文件。其任务即在于定量刻画预热平衡阶段的升温过程。建模时取药材内部一薄圆筒为控制体，分别对其建立能量守恒与质量守恒关系，导出以温度和水分浓度为未知量的导热方程与扩散方程；再补充轴心处的对称性条件、表面处的对流换热与表面蒸发条件，以及入烘房时均匀的初始状态，构成封闭的定解问题，得到前1800\,s药材内部的温度与水分的空间分布。

\subsection*{问题二}
问题二要求求解恒温干燥阶段时，各个时间点距药材中心不同位置的温度和水分浓度。物性改用附录 3 中随温度与含水率变化的表达式。首先，根据附件 1，拟合出恒温干燥阶段烘房的温度为 50◦C 和含水量为 0.05kg/kg。其物理框架与问题一一致，仍以薄圆筒的守恒关系建立方程，边界条件不变，初值取问题一终止时刻的分布。区别在于物性不再取常数：扩散系数依赖于温度，而蒸发耗热又取决于水分到达表面的速率，温度场与水分场经由物性相互耦合，联立求解，得到较高精度结果。

\subsection*{问题三}
问题三沿用问题二的物性与几何设定，要求计算药材内部含水率最大值降至 0.15 kg/kg 所需的持续时间，同时要求给出该过程中的温度和含水量的分布。首先，判断“干燥终点”位置为圆柱中心轴线，并予以证明；然后，重心由时间推进转为终止判据的确定，需构造以当前最高含水率与阈值之差为取值的监测函数，由其符号变化确定过程终止的时刻；最后，采用连续事件定位与空间插值实现数值求解与结果提取，得到干燥停止的时间。

\subsection*{问题四}
问题四要求在问题三基础上进行失水收缩分析，确定药材的烘干时长，给出水分浓度随时间、空间变化的结果。首先，建立药材的移动边界模型，物性改用附录4。然后，通过坐标变换推导材料坐标系下的控制方程与表面缩放边界，将移动域转化为固定域进行计算。此时，原问题即化归为与前三问同构的形式；作此变换后药材内部质量守恒与温度通量计算需要在求解过程中同步进行映射。

\begin{figure}[H]\centering
\begin{tikzpicture}[node distance=8mm and 11mm]
\node[box] (a) {附件与物性\\环境、半径、经验式};
\node[box,right=of a] (b) {局部守恒建模\\通量、初值、边界};
\node[box,right=of b] (c) {有限体积离散\\稀疏隐式积分};
\node[box,below=of c] (d) {分布与阈值事件\\四问结果输出};
\node[box,left=of d] (e) {独立解与收敛\\守恒及敏感性};
\node[box,left=of e] (f) {条件化结论\\工艺解释与改进};
\draw[arr](a)--(b);\draw[arr](b)--(c);\draw[arr](c)--(d);
\draw[arr](d)--(e);\draw[arr](e)--(f);
\end{tikzpicture}
\caption{建模、求解与验证的总体技术路线}\label{fig:roadmap}
\end{figure}

\nextpage
\section{模型假设}
为使给定数据能够闭合为可计算问题，并保持论文与既有算法一致，采用以下假设。

\begin{enumerate}
\item \textbf{圆柱径向主导。}药材形状近似为圆柱，材料及外部交换在周向上均匀。端面总面积与侧面积之比为$R/L=0.08$，这一几何关系支持忽略端面与空气的传热和传质。
\item \textbf{显热近似。}忽略水分蒸发的潜热。即圆柱侧面的水分蒸发不降低本体温度。
\item \textbf{有效湿度边界。}干燥理论告诉我们对流传质的驱动力由药材的当前含水率与干基含水率之差决定。本题没有给平衡含水率，我们使用空气水分浓度做近似。
\item \textbf{环境与交换条件。}两个对流系数主要和空气流速相关，其次是温度和相对湿度，以及药材的形状。为了简化问题，方便计算，我们认为未给出$h$和$h_m$的后三问的空气流速条件与第一问基本一致，即两个对流交换系数与第一问取值相同。
\item \textbf{均匀径向收缩。}第四问中长度固定，外半径由附件2线性插值；内部骨架速度取$v=r\dot R/R$，不考虑开裂、分层和不均匀变形。固定$x=r/R(t)$对应同一材料位置。
\end{enumerate}

\nextpage
\section{符号说明}
\begin{table}[H]\centering
\caption{主要符号及其单位}\label{tab:symbols}
\small
\begin{tabularx}{\linewidth}{>{\centering\arraybackslash}p{24mm}X>{\centering\arraybackslash}p{34mm}}
\toprule
符号 & 说明 & 单位\\\midrule
$r,\ R(t),\ L$ & 径向坐标、外半径、圆柱长度 & m\\
$t,\ t_*$ & 时间、干燥临界时间 & s\\
$T,\ T_a,\ T_s$ & 药材、空气及表面温度 & $^\circ$C\\
$T_K$ & 计算经验扩散系数使用的绝对温度 & K\\
$C,\ C_s,\ C_e$ & 内部、表面及有效平衡干基含水率 & kg/kg\\
$m_d,\ m_w$ & 干物质量、水质量 & kg\\
$\rho,\ \rho_d$ & 有效热密度、单位体积参考干物质量 & kg/m$^3$\\
$c_p$ & 有效比热容 & J/(kg$\cdot$K)\\
$k$ & 导热系数 & W/(m$\cdot$K)\\
$D$ & 有效水分扩散系数 & m$^2$/s\\
$h$ & 表面对流换热系数 & W/(m$^2\cdot$K)\\
$h_m$ & 表面有效传质系数 & m/s\\
$q,\ j$ & 向外热流密度、相对骨架水质量通量 & W/m$^2$；kg/(m$^2\cdot$s)\\
$v$ & 径向骨架速度 & m/s\\
$x,\ x_i$ & 归一化坐标及其节点位置 & 1\\
$N,\ \beta$ & 径向网格区间数、加密参数 & 1\\
$\omega_i$ & 归一化控制体积的径向权重 & 1\\
$H_{i+1/2}$ & 归一化热通量，定义见式\eqref{eq:fluxes} & W/m$^3$\\
$W_{i+1/2}$ & 归一化含水率通量 & s$^{-1}$\\
$\alpha,\ Bi$ & 热扩散率、基于半径的热毕渥数 & m$^2$/s；1\\
$M(t),\ \overline C(t)$ & 全域最大、干质量加权平均含水率 & kg/kg\\
$\bm y,\ \bm F$ & 离散状态向量、其时间变化率 & 分量单位各异\\
\bottomrule
\end{tabularx}
\end{table}
